\section{Kết luận}

Nhóm đã xây dựng một hệ thống hoàn chỉnh cho Đề tài 4: từ dữ liệu bản đồ thật, qua các cấu trúc dữ liệu tự cài đặt, tới giao
diện demo cho phép tìm tuyến, gây sự cố giao thông, xem mạng xương sống giữa các kho, phân cụm đơn hàng và lập chuyến xe.
Lõi giải thuật khoảng 1\,300 dòng Python không dùng \texttt{dict}, \texttt{set} hay \texttt{heapq}, và quy tắc này được một
test phân tích cú pháp kiểm tra tự động. Mỗi giải thuật có test đối chiếu với lời giải chuẩn: Dijkstra với Bellman-Ford,
Kruskal với Prim, Knapsack và Held-Karp với vét cạn, Union-Find với BFS.

Kết quả thực nghiệm khớp với phân tích Big-O ở những chỗ mà lý thuyết dự đoán xu hướng: cận $O(\alpha(n))$ của Union-Find,
độ tăng $O(n^2 2^n)$ so với $O(n!)$ của TSP, tính tuyến tính theo $n$ của Knapsack khi $W$ cố định. Thực nghiệm cũng chỉ ra hai
chỗ cần đọc lý thuyết cẩn thận hơn. Thứ nhất, lợi thế của Adaptable Priority Queue phụ thuộc vào tỉ lệ decrease-key, vốn rất
thấp trên mạng đường. Thứ hai, ở quy mô vừa phải, hằng số của trình thông dịch có thể làm đảo thứ hạng giữa hai giải thuật có
cận khác nhau, như trường hợp Prim và Kruskal.

\textbf{Hướng phát triển.} Dùng A* với heuristic khoảng cách đường chim bay và Dijkstra hai chiều để giảm số đỉnh phải duyệt;
cập nhật đường đi ngắn nhất tăng dần thay vì tính lại; rút gọn các đỉnh bậc 2 của dữ liệu OSM; mở rộng bài toán lập chuyến
thành VRP nhiều xe có khung giờ giao; cài lại phần lõi bằng C++ để so sánh các giải thuật mà không bị ảnh hưởng bởi hằng số
của trình thông dịch.
